\chapter{Técnicas Digitales II: Control Embebido, Arquitectura RT y HMI}
\label{sec:tecnicas_digitales_2}

En el área de Técnicas Digitales II, el anteproyecto abarca la concepción de la arquitectura del sistema embebido de
control, la adopción de un esquema de ejecución multitarea en tiempo real (RTOS), la interpretación de comandos de
mecanizado estándar y la gestión integral de periféricos e interfaz de usuario por medio de la biblioteca CMSIS.

\section{Arquitectura del Sistema Embebido y Ejecución Multitarea}
\label{sec:td2_freertos}

Para responder de forma determinística ante eventos críticos de seguridad (como la detección de tope o sobrecorriente)
y sostener simultáneamente una generación fluida y precisa de pulsos de movimiento sin bloquear la interacción con el
usuario, el software se estructura bajo un sistema operativo en tiempo real (RTOS).

La arquitectura se organiza en capas de procesamiento desacopladas, asignando prioridades de acuerdo con la
criticidad temporal de cada función:

\begin{itemize}
  \item \textbf{Capa de Seguridad y Eventos Críticos (Prioridad Máxima):} Permanece a la espera de señales o
    interrupciones de protección provenientes de la etapa analógica o pulsadores de parada. Ante un evento, toma el
    control inmediato para inhibir los trenes de pulsos y proteger la integridad mecánica.
  \item \textbf{Capa de Control Cinemático y Generación de Movimiento (Prioridad Alta):} Procesa bloques de trayectoria
    desde colas de memoria, calcula perfiles de aceleración/desaceleración y comanda temporizadores por hardware para
    la emisión precisa de señales de avance (paso y dirección).
  \item \textbf{Capa de Interpretación de Comandos (Prioridad Media):} Analiza las cadenas de caracteres de mecanizado
    (código G), valida su sintaxis y convierte las coordenadas geométricas en desplazamientos escalados para el
    planificador de movimiento.
  \item \textbf{Capa de Almacenamiento y Comunicaciones (Prioridad Media-Baja):} Administra la transferencia de
    archivos de trabajo desde memorias extraíbles (ej. tarjeta SD) o enlaces serie con una estación de control superior.
  \item \textbf{Capa de Interfaz Hombre-Máquina (Prioridad Baja):} Gestiona la visualización del estado del sistema en
    pantallas locales y lee controles de entrada de usuario (teclados, encoders o pulsadores) para funciones auxiliares
    y posicionamiento manual.
\end{itemize}

En la figura~\ref{fig.ej3:arquitectura_freertos} se ilustra la organización conceptual de las tareas concurrentes y los
mecanismos de comunicación interproceso tales como colas, semáforos y grupos de eventos.

\begin{figure}[!ht]
  \centering
  \resizebox{0.88\linewidth}{!}{\begin{tikzpicture}[
  >=stealth, thick,
  font=\small,
  node distance=1.1cm and 1.6cm,
  task/.style={draw=blue!70!black, fill=blue!10, rectangle, rounded corners=4pt, minimum height=1.0cm, minimum width=2.8cm, align=center},
  queue/.style={draw=yellow!60!black, fill=yellow!20, ellipse, minimum height=0.8cm, minimum width=2.4cm, align=center, font=\footnotesize},
  safety/.style={draw=red!70!black, fill=red!12, rectangle, rounded corners=4pt, minimum height=1.0cm, minimum width=2.8cm, align=center, font=\small\bfseries},
  isr/.style={draw=red!80!black, fill=red!20, rectangle, minimum height=0.8cm, minimum width=2.4cm, align=center, font=\footnotesize\bfseries}
]

  % Columna Central: Pipeline de Ejecución de Mecanizado
  \node[task] (storage) {Entrada de Datos\\{\footnotesize (Lectura SD / UART)}};
  \node[queue, below=0.9cm of storage] (qcode) {Cola de Comandos\\{\footnotesize (\texttt{Queue G-Code})}};
  \node[task, below=0.9cm of qcode] (parser) {Intérprete G-Code\\{\footnotesize (Validación y Conversión)}};
  \node[queue, below=0.9cm of parser] (qmotion) {Cola de Trayectoria\\{\footnotesize (\texttt{Queue Motion})}};
  \node[task, below=0.9cm of qmotion] (planner) {Planificador de Movimiento\\{\footnotesize (Generación STEP/DIR)}};

  % Columna Izquierda: Interfaz de Usuario
  \node[task, left=of parser] (hmi) {Gestión HMI\\{\footnotesize (Pantalla / Encoder)}};

  % Columna Derecha: Sistema de Seguridad
  \node[isr, right=of storage] (exti) {Evento EXTI\\{\footnotesize (Tope Analógico)}};
  \node[safety, right=of qcode] (safetytask) {Supervisión\\y Seguridad};

  % Conexiones del Pipeline Central
  \draw[->, draw=blue!80!black] (storage) -- (qcode);
  \draw[->, draw=blue!80!black] (qcode) -- (parser);
  \draw[->, draw=blue!80!black] (parser) -- (qmotion);
  \draw[->, draw=blue!80!black] (qmotion) -- (planner);

  % Conexión HMI con Planificador y Parser
  \draw[<->, draw=cyan!80!black] (hmi) |- (planner);
  \draw[->, draw=cyan!80!black] (hmi) -- (parser);

  % Conexiones de Seguridad
  \draw[->, draw=red!80!black] (exti) -- node[midway, fill=white, inner sep=2pt, font=\footnotesize] {Semáforo} (safetytask);
  \draw[->, dashed, draw=red!80!black] (safetytask) |- node[pos=0.7, fill=white, inner sep=2pt, font=\footnotesize] {Inhibición Inmediata} (planner);

\end{tikzpicture}
}
  \caption{Esquema de arquitectura de software en tiempo real con separación por capas de prioridad y mecanismos de sincronización.}
  \label{fig.ej3:arquitectura_freertos}
\end{figure}

\section{Módulo Intérprete de Código G y Perfilado de Movimiento}
\label{sec:td2_gcode}

El sistema embebido incorporará un módulo analizador sintáctico para interpretar instrucciones normalizadas según el
estándar ISO 6983 / RS-274D. La implementación se enfocará en un subconjunto fundamental que comprende:
\begin{itemize}
  \item Movimientos de posicionamiento rápido no interpolado.
  \item Interpolación lineal con velocidad de avance controlada (\textit{feedrate}).
  \item Comandos de temporización y pausas programadas.
  \item Comandos auxiliares de control de herramientas (encendido, apagado y ajuste de velocidad del \textit{spindle}).
\end{itemize}

En la figura~\ref{fig.ej3:flujo_gcode} se describe el flujo general de datos desde la recepción de la instrucción hasta
su ejecución en los actuadores.

\begin{figure}[!ht]
  \centering
  \begin{tikzpicture}[node distance=1.5cm, auto, >=stealth, thick,
  startstop/.style={draw, fill=green!15, rectangle, rounded corners, minimum height=0.8cm, text centered, align=center},
  process/.style={draw, fill=blue!10, rectangle, minimum height=1.0cm, text centered, align=center},
  decision/.style={draw, fill=yellow!15, diamond, aspect=2, inner sep=2pt, text centered, align=center}]

  \node [startstop] (start) {Recepción Cadena G-Code};
  \node [process, below=of start] (parse) {Extracción de Comando y Parámetros (X, Y, Z, F, S)};
  \node [decision, below=1.2cm of parse] (valid) {¿Sintaxis Válida?};
  \node [process, right=2.0cm of valid] (error) {Enviar Error por HMI / UART};
  \node [process, below=1.4cm of valid] (calc) {Transformación a Pasos ($N_x, N_y, N_z$)};
  \node [process, below=of calc] (plan) {Perfilado de Velocidad (Rampa Trapezoidal)};
  \node [process, below=of plan] (enqueue) {Insertar en Cola de Movimiento (\texttt{qMotion})};
  \node [startstop, below=of enqueue] (exec) {Timer genera pulsos STEP/DIR};

  \draw [->] (start) -- (parse);
  \draw [->] (parse) -- (valid);
  \draw [->] (valid) -- node[above, font=\small] {No} (error);
  \draw [->] (valid) -- node[right, font=\small] {Sí} (calc);
  \draw [->] (calc) -- (plan);
  \draw [->] (plan) -- (enqueue);
  \draw [->] (enqueue) -- (exec);
  \draw [->] (error.north) |- (start.east);
\end{tikzpicture}

  \caption{Diagrama de flujo para el procesamiento, validación y ejecución de comandos de mecanizado.}
  \label{fig.ej3:flujo_gcode}
\end{figure}

\section{Selección del Microcontrolador y Asignación de Periféricos}
\label{sec:td2_perifericos}

Como unidad central de procesamiento se utilizará un microcontrolador reciclado \textbf{GD32F103RCT6} (arquitectura ARM Cortex-M3), un reemplazo directo compatible con el STM32F103RCT6 capaz de operar a una frecuencia máxima de reloj de \SI{108}{\mega\hertz}. Al pertenecer a la misma familia de microcontroladores que la plataforma base recomendada por la cátedra (Blue Pill / STM32F103C8T6), se garantiza plena compatibilidad a nivel de arquitectura, periféricos, bibliotecas CMSIS y soporte del kernel FreeRTOS. Adicionalmente, la adopción de este componente reutilizado en encapsulado LQFP64 responde a la necesidad de contar con un adecuado margen de reserva (\textit{headroom}) de líneas de entrada/salida y temporizadores por hardware, previniendo restricciones de conectividad al definir la cantidad final de pines requeridos para los controladores de movimiento (motores paso a paso), el accionamiento del \textit{spindle}, los canales de sensado analógico y la interfaz HMI.

Para satisfacer las demandas del sistema, el dispositivo asignará el siguiente conjunto de capacidades periféricas:

\begin{itemize}
  \item \textbf{Temporizadores por Hardware (Timers):} Módulos de temporización con capacidad de interrupción periódica
    y modulación PWM para la generación determinística de frecuencias de paso y control del \textit{spindle} sin sobrecargar la CPU.
  \item \textbf{Buses de Comunicación Serie:} Interfaces sincrónicas de alta velocidad (como SPI/SDIO) para el acceso al
    almacenamiento masivo (tarjeta SD) e interfaces serie (como UART o I2C) para comunicación y periféricos de HMI.
  \item \textbf{Canales de Conversión A/D (ADC):} Lectura periódica de las tensiones provenientes del acondicionamiento analógico
    para el seguimiento del consumo energético, estado de carga y sensado de corriente.
  \item \textbf{Líneas GPIO e Interrupciones Externas (EXTI):} Control de señales de habilitación de drivers y captura
    inmediata de eventos asíncronos de seguridad.
\end{itemize}
